The theory predicts when a balanced affinity helps. What matters is how a
velocity field splits into a gradient part and a circulating part. We
decompose the edge flow $f_{ij}=\tfrac12(v_i+v_j)^\top(x_j-x_i)$ on a
30-nearest-neighbour graph in 30 principal components, following
\citet{jiang2011statistical}. The null model permutes the centred
velocities among cells.

\begin{table}[t]
  \caption{Hodge content of real RNA-velocity fields: share of
    $\|f\|^2$ carried by the gradient part, its permutation null, and the
    part that carries the ground-truth progress. Gastrulation:
    $n=30{,}000$ cells, E6.5--E8.5 \citep{pijuansala2019single}. RPE1:
    $n=2{,}793$ cells with FUCCI phase \citep{battich2020sequencing}, velocity
    from dynamo \citep{qiu2022mapping}.}
  \label{tab:hodge}
  \centering\small
  \begin{tabular}{@{}lccl@{}}
    \toprule
    velocity & gradient & null & progress carried by \\
    \midrule
    gastrulation, direct     & 0.93 & 0.33 & gradient \\
    gastrulation, scVelo     & 0.78 & 0.31 & gradient \\
    RPE1, direct             & 0.67 & 0.21 & non-gradient \\
    RPE1, dynamo kernel      & 0.29 & 0.24 & mostly non-gradient \\
    \bottomrule
  \end{tabular}
\end{table}

\Cref{tab:hodge} shows two opposite regimes.
\begin{itemize}
  \item \textbf{Developmental trajectories are gradient-dominated.} On
    gastrulation, a walk along the gradient part alone recovers the whole
    forward-in-time signal. The remainder carries none of it. A large share
    of the velocity is also a vector common to all cells, which
    \cref{cor:gauge} makes invisible to any balanced affinity.
  \item \textbf{The cell cycle circulates.} On RPE1, the gradient part of
    the direct velocity makes walks go slightly \emph{backwards} in phase.
    The non-gradient part carries the progress around the cycle; for the
    kernel velocity it carries most of it over long walks. The
    velocity circulates around closed loops through the phase bins far
    above the null ($z$ between $5$ and $18$).
\end{itemize}

\paragraph{Balanced versus row-stochastic affinities.}
\Cref{tab:dsrs} builds both affinities from the same velocity, with the
per-row perplexity calibration of practical neighbour embeddings.
\begin{table}[t]
  \caption{Balanced (Sk) versus row-stochastic (RS) Randers affinities
    built from the same velocity: squared cost, per-row perplexity
    $\sqrt n$, dense. ``cos'' is the mean cosine between each cell's shift
    in outgoing displacement, relative to the symmetric affinity, and the
    gradient or non-gradient Hodge part of the velocity. ``progress'' is
    the ground-truth progress after $10$ steps of the walk, minus that of
    the symmetric walk: pseudotime units for pancreas, fractions of a cycle
    for RPE1. The linear cost gives the same picture.}
  \label{tab:dsrs}
  \centering\small
  \begin{tabular}{@{}llrrr@{}}
    \toprule
    data & norm. & cos grad. & cos rest & progress \\
    \midrule
    pancreas      & RS & $0.56$ & $0.21$ & $0.167$ \\
                  & Sk & $0.38$ & $0.26$ & $0.000$ \\
    \addlinespace
    RPE1, direct  & RS & $0.56$ & $0.14$ & $0.060$ \\
                  & Sk & $0.01$ & $0.33$ & $0.082$ \\
    \addlinespace
    RPE1, kernel  & RS & $0.34$ & $0.36$ & $0.311$ \\
                  & Sk & $0.12$ & $0.37$ & $0.241$ \\
    \bottomrule
  \end{tabular}
\end{table}

Three observations follow the theory.
\begin{itemize}
  \item \textbf{No net progress along a potential.} On pancreas, the
    balanced walk's mean pseudotime gain is zero to numerical precision,
    as \cref{prop:netflow} requires. The row-stochastic walk gains steadily.
  \item \textbf{Circulation survives balancing.} On RPE1, the balanced walk
    keeps the progress around the cycle. With the direct velocity it
    progresses further than the row-stochastic walk, whose extra
    sensitivity to the gradient part pushes it backwards in phase.
  \item \textbf{The balanced shift drops the gradient part.} On RPE1 its
    correlation with the gradient part falls to about zero for the direct
    velocity, while its correlation with the remainder rises. On pancreas
    the balanced shift still has a sizeable gradient correlation. This is
    the first-order leak of per-row bandwidths (\cref{cor:leak}). By
    \cref{prop:netflow} it must be compensated, and it is: per pseudotime
    decile, the balanced chain moves cells forward in the first seven and
    backwards in the last three. Its largest step, $-0.076$, is in the
    terminal decile. The row-stochastic chain moves forward in every
    decile.
\end{itemize}


\paragraph{Consequence for embeddings.}
These numbers explain an observation from fitting Randers neighbour
embeddings on pancreas \citep{bastidasponce2019comprehensive}. With a
row-stochastic input affinity, the fitted flow agreed with the true
velocity at cosine $0.80$ and with scVelo's projection at $0.76$. With a
balanced one, these fell to $0.69$ and $0.54$. The balanced direction was
attenuated rather than erased, as \cref{cor:leak} predicts: per-row
bandwidths and the squared cost leak part of the gradient drift back.
